\title{xLSTM: Extended Long Short-Term Memory}

\begin{abstract}
The fundamental principles behind Long Short-Term Memory (LSTM) networks—namely, the constant error carousel and gating—were established during the 1990s. Over the decades, LSTMs have proven remarkably durable, significantly shaping the development of deep learning methods, and notably serving as the foundation for the first generation of Large Language Models (LLMs). Nonetheless, the emergence of the Transformer, powered by highly parallelizable self-attention mechanisms, signaled a transformative shift and ultimately surpassed LSTMs when applied at scale. This work revisits LSTMs by asking: What performance can be achieved in language modeling if LSTMs are scaled up to billions of parameters, utilize state-of-the-art techniques employed in contemporary LLMs, and address documented weaknesses intrinsic to LSTM models? To that end, we first propose exponential gating, complemented by suitable normalization and stabilization strategies. Next, we alter the core LSTM memory mechanism, resulting in two variants: (i) sLSTM, which leverages a scalar memory representation with scalar updates and introduces a revised method for memory mixing, and (ii) mLSTM, which is entirely parallelizable and features a matrix-based memory alongside a covariance-inspired update process. By incorporating these advanced LSTM variants within residual block frameworks, we construct xLSTM blocks, which are subsequently organized into larger xLSTM architectures through residual stacking. The combination of exponential gating and enhanced memory structures equips xLSTM models to rival, or even exceed, the performance and scaling attributes of the most advanced Transformers and State Space Models.\\
Code available at: \url{https://github.com/NX-AI/xlstm}
\end{abstract}

\section{Introduction}
\label{sec:intro}

The Long Short-Term Memory (LSTM) ideas~\citep{Hochreiter:91,Hochreiter:97nips,Hochreiter:97}, i.e., the constant error carousel and gating, were introduced to overcome the vanishing gradient problem of recurrent neural networks~\citep{Hochreiter:91,Hochreiter:00book}:

\begin{align}
\label{eq:lstm_idea}
  c_t \ = \ f_t \ c_{t-1} \ + \ 
          i_t \  z_t \ , \quad  
          h_t \ = \ o_t \ \psi({c_t}) \ . 
\end{align}

The constant error carousel refers to the process where the previous cell state $c_{t-1}$ (depicted in green) is incrementally updated by the incoming inputs $z_t$, with sigmoid gates (in blue) regulating this addition. The input gate $i_t$ together with the forget gate $f_t$ modulate how the cell state evolves, whereas the output gate $o_t$ determines what portion of the memory cell becomes the hidden state $h_t$. To obtain the hidden state, the cell state is first processed through the nonlinear function $\psi$ before being filtered by the output gate.

Long Short-Term Memory networks (LSTMs) have achieved notable success across a wide array of domains \citep{Hochreiter:01,Hochreiter:07,Schmidhuber:15}, dominating text generation tasks until the emergence of Transformers in 2017~\citep{Vaswani:17}. Their performance has been substantiated by impressive results in multiple sequential data processing challenges, including but not limited to text generation~\citep{Graves:13,Karpathy:15blog}, handwritten sequence synthesis~\citep{Graves:13}, sequence-to-sequence translation~\citep{Sutskever:14nips}, evaluation of computer code~\citep{Zaremba:14arxiv}, captioning images~\citep{Karpathy:15,Hossain:19}, automatic source code production~\citep{Karpathy:15blog}, modeling rainfall-runoff processes~\citep{Kratzert:18,Kratzert:19}, and developing hydrological models for flood forecasting~\citep{Nearing:24}. Within the field of reinforcement learning, LSTM-based architectures remain the most effective sequence models, as demonstrated by their use in prominent systems such as AlphaStar for StarCraft II~\citep{Vinyals:17short}, OpenAI Five for Dota 2~\citep{OpenAI:19}, and magnetic control frameworks for nuclear fusion reactors~\citep{Degrave:22short}. A key strength of LSTMs lies in their capacity to form abstractions, effectively capturing and encoding semantic information in their internal memory structures~\citep{Karpathy:15blog}. This ability has been illustrated by the discovery of distinct neural units associated with numerical reasoning and syntax~\citep{Lakretz:19}, linguistic features~\citep{Bau:19}, and sentiment representation~\citep{Radford:17arxiv}. Despite advances in other architectures, LSTMs continue to serve important roles in cutting-edge applications~\citep{Degrave:22short,Nearing:24}, attesting to their enduring utility.

While LSTMs have achieved notable achievements, they are hindered by three key drawbacks: (i) Difficulty in revising storage choices. This shortcoming is illustrated with the {\it Nearest Neighbor Search} task (refer also to Appendix~\ref{sec:appNearestNeighborSearch}): given a reference vector, the sequence must be scanned one element at a time to find and retrieve the most similar vector's associated value at the sequence's end. The leftmost panel of Figure~\ref{fig:lstmProblems} displays mean squared errors for this challenge. Standard LSTM architectures find it hard to update a value once stored if a closer match appears later in the sequence; in contrast, our xLSTM addresses this via the introduction of exponential gating. (ii) Restricted capacity for information storage, as all data must be condensed into scalar-valued cell states. This is demonstrated using the {\it Rare Token Prediction} task. In Figure~\ref{fig:lstmProblems} (right panel), perplexities for token prediction on Wikitext-103~\citep{Merity:17} are reported across frequency bands. LSTM models perform suboptimally for infrequent tokens, reflecting their constrained storage capability. The xLSTM overcomes this constraint through matrix-based memory storage. (iii) Inherent sequential processing restrictions resulting from the way memory content is blended via hidden-to-hidden transitions, which severely limits parallelization opportunities. 

Such shortcomings have been instrumental in motivating the development and widespread adoption of Transformer models~\citep{Vaswani:17} for language modeling tasks. However, an open question remains: if we can remove these limitations and scale LSTM variants, such as our xLSTM, to match the scale of contemporary Large Language Models, what levels of performance in language modeling might we achieve?

\section{Extended Long Short-Term Memory}
\label{sec:xLSTM}

To address the constraints inherent in LSTM architectures, Extended Long Short-Term Memory (xLSTM) presents two significant enhancements over the foundational LSTM concept from Equation~\eqref{eq:lstm_idea}. These enhancements—namely, exponential gating and innovative memory constructs—expand the LSTM family to include two additional variants: (i) the newly proposed sLSTM (refer to Section~\ref{sec:sLSTM}), characterized by its scalar memory, scalar update mechanism, and incorporation of memory mixing, and (ii) the mLSTM (see Section~\ref{sec:mLSTM}), which employs a matrix-based memory alongside a covariance (outer product) update rule that allows for complete parallelism. Both sLSTM and mLSTM introduce exponential gating as a means of improving LSTM capabilities. In pursuit of parallel execution, mLSTM forgoes memory mixing, specifically the recurrent connections among hidden states. The architecture of both sLSTM and mLSTM can be expanded to accommodate multiple memory cells; notably, sLSTM permits mixing of memory across its cells. Additionally, sLSTM supports configurations with multiple heads—where memory mixing occurs among cells within the same head, but not across distinct heads—yielding a novel mode of memory mixing when combined with exponential gating. In the context of mLSTM, adopting multiple heads is functionally equivalent to expanding the number of memory cells.

When these variants are embedded within residual block modules, the resultant constructs are termed xLSTM blocks (outlined in Section~\ref{sec:xLSTMblock}). Layering these xLSTM blocks residually enables the creation of full xLSTM architectures (discussed in Section~\ref{sec:xLSTMarchitecure}). A visual schematic of the xLSTM architecture and its constituent elements is depicted in Figure~\ref{fig:xlstm_sketch}.

\subsection{Review of the Long Short-Term Memory}
\label{sec:orgLSTM}

The initial concept of LSTM~\citep{Hochreiter:91,Hochreiter:97nips,Hochreiter:97} featured a scalar memory cell serving as the main unit for computation and information storage, specifically designed to mitigate the vanishing gradient problem~\citep{Hochreiter:91,Hochreiter:00book} by employing the constant error carousel, which is implemented via updates to the cell state. This memory cell architecture utilizes three gating mechanisms: the input, output, and forget gates, with the forget gate having been subsequently introduced by \citet{Gers:00}. The equations governing the LSTM memory cell dynamics at each time point $t$ are as follows:

\begin{align}
c_t \ &= \  f_t \ c_{t-1} \ + \ i_t \ z_t &  & &\text{cell state} \\
h_t  \ &= \ o_t \ \tilde{h}_t \ , & \tilde{h}_t \ &= \ \psi \left( c_t\right)
  &\text{hidden state} \\
z_t \ &= \ \varphi \left( \tilde{z}_t \right) \ , 
  &\tilde{z}_t \ &=  \ \mathbf{w}^\top_{z} \ \mathbf{x}_t \ + \
  r_{z}  h_{t-1} \ + \  b_{z} \ \
  &\text{cell input} \\
i_t \ &= \ \sigma \left( \tilde{i}_t  \right) \ , 
  &\tilde{i}_t \ &= \ \mathbf{w}^\top_{i} \ \mathbf{x}_t \ + \
  r_{i}  \ h_{t-1} \ + \  b_{i} \ \
  &\text{input gate} \\
f_t \ &= \ \sigma \left(  \tilde{f}_t \right) \ , 
  &\tilde{f}_t \ &= \ \mathbf{w}^\top_{f} \ \mathbf{x}_t  \ + \
  r_{f}  \ h_{t-1} \ + \  b_{f} \ \
  &\text{forget gate} \\
o_t \ &= \ \sigma \left( \tilde{o}_t \right) \ , 
  &\tilde{o}_t  \ &= \ \mathbf{w}^\top_{o} \ \mathbf{x}_t \ + \
  r_{o}  \ h_{t-1} \ + \  b_{o} \ \
  &\text{output gate} 
\end{align}

The vectors $\mathbf{w}_{z}$, $\mathbf{w}_{i}$, $\mathbf{w}_{f}$, and $\mathbf{w}_{o}$ denote the input weight vectors linking the input $\mathbf{x}_t$ to the cell input, input gate, forget gate, and output gate, respectively. The parameters $r_{z}$, $r_{i}$, $r_{f}$, and $r_{o}$ represent the recurrent weights connecting the previous hidden state $h_{t-1}$ to the corresponding cell input, input gate, forget gate, and output gate. The terms $b_{z}$, $b_{i}$, $b_{f}$, and $b_{o}$ serve as the respective biases. The cell input and hidden state activation functions are given by $\varphi$ and $\psi$, both commonly instantiated as $\tanh$. The activation function $\psi$ is responsible for ensuring the cell state remains within bounds, as otherwise it may grow without restriction. All gating mechanisms employ the sigmoid activation, specified as $\sigma \left( x \right)= 1/(1+\exp(-x))$. Later versions introduced the combination of several scalar memory cells $\mathbf{c}_t \in \mathbb{R}^{d}$ into a single vector, facilitating the application of recurrent weight matrices $\mathbf{R} \in \mathbb{R}^{d \times d}$, which enables mixing among the outputs of different memory cells~\citep{Greff:15}; further information is provided in Appendix~\ref{sec:apporgLSTM}. Ablation experiments demonstrated that each element within the memory cell architecture is essential~\citep{Greff:15}.

\subsection{sLSTM}
\label{sec:sLSTM}

To enhance LSTMs' capacity to modify how information is retained, we incorporate exponential gates (depicted in red), as well as normalization and stabilization mechanisms. Specifically, we allow the input and forget gates to utilize exponential activation functions. For normalization purposes, we implement a normalizer state, which accumulates the sum of the input gate multiplied by all subsequent forget gates.

The scalar sLSTM forward pass is:

\begin{align}
\label{eq:slstmforward}
c_t \ &= \  f_t \ c_{t-1} \ + \ i_t \ z_t & & 
 &\text{cell state} \\
n_t \ &= \  f_t \ n_{t-1} \ + \ i_t  & & 
 &\text{normalizer state} \\
h_t \ &= \ o_t \ \tilde{h}_t \ ,
  &\tilde{h}_t \ &= \ c_t / n_t
&\text{hidden state} \\
z_t \ &= \ \varphi \left( \tilde{z}_t \right) \ , 
  &\tilde{z}_t \ &=  \ \mathbf{w}^\top_{z} \ \mathbf{x}_t \ + \
  r_{z}  h_{t-1} \ + \  b_{z}
  &\text{cell input} \\
i_t \ &= \ \!\exp \left( \tilde{i}_t  \right) \ , 
  &\tilde{i}_t \ &= \ \mathbf{w}^\top_{i} \ \mathbf{x}_t \ + \
  r_{i}  \ h_{t-1} \ + \  b_{i}
  &\text{input gate} \\
f_t \ &= \ \sigma \left(  \tilde{f}_t \right) \ \text{OR} \
\!\exp \left(  \tilde{f}_t \right) \ ,
  &\tilde{f}_t \ &= \ \mathbf{w}^\top_{f} \ \mathbf{x}_t  \ + \
  r_{f}  \ h_{t-1} \ + \  b_{f}
  &\text{forget gate} \\
o_t \ &= \ \sigma \left( \tilde{o}_t \right) \ , 
  &\tilde{o}_t  \ &= \ \mathbf{w}^\top_{o} \ \mathbf{x}_t \ + \
  r_{o}  \ h_{t-1} \ + \  b_{o}
  &\text{output gate} 
\end{align}

We transfer the original LSTM gating techniques, i.e., input- and/or hidden-dependent gating plus bias term, to the new architectures. Exponential activation functions can lead to large values that cause overflows. Therefore, we stabilize gates with an additional state \mbox{$m_t$~\citep{Milakov:18arxiv}}:

\begin{align}
\label{eq:slstmstabil}
m_t \ &= \ \max \left( \log ( f_t ) + m_{t-1} , \log ( i_t ) \right) &\text{stabilizer state} \\
i'_t \ &= \ \!\exp \left( \log \left ( i_t \right) - m_t \right) \ = \ \!\exp \left( \tilde{i}_t - m_t \right) \ \, 
  &\text{stabil. input gate} \\
f'_t \ &= \ \!\exp \left( \log \left( f_t \right) + m_{t-1} - m_t \right) \ \, 
  &\text{stabil. forget gate}
\end{align}

Appendix~\ref{sec:appsLSTM} demonstrates that substituting $f_t$ with $f'_t$ and $i_t$ with $i'_t$ during the forward computation leaves both the network’s output and the gradients of the loss with respect to the model’s parameters unchanged.

\paragraph{New Memory Mixing.} Like standard LSTM architectures, sLSTM may comprise several memory cells (refer to Appendix~\ref{sec:appsLSTM}). The presence of multiple memory cells allows for mixing of information through recurrent connections $\mathbf{R}_{\mathbf{z}}$, $\mathbf{R}_{\mathbf{i}}$, $\mathbf{R}_{\mathbf{f}}$, and $\mathbf{R}_{\mathbf{o}}$, which link the hidden state vector $\mathbf{h}$ with the memory cell input $\mathbf{z}$ and the gates $\mathbf{i}$, $\mathbf{f}$, and $\mathbf{o}$. Exponential gating introduces a distinct feature in the process of memory mixing. In the new sLSTM framework, several heads are permitted, facilitating memory mixing within each head but not across different heads. The combination of the head structure in sLSTM and exponential gating results in a novel mechanism for mixing memories.

\subsection{mLSTM}
\label{sec:mLSTM}

To expand the storage capacity of LSTMs, the memory cell is generalized from a single scalar value $c \in \mathbb{R}$ to a matrix representation $\mathbf{C} \in \mathbb{R}^{d \times d}$. As a consequence, information retrieval is achieved through matrix multiplication operations. During each timestep $t$, we store two vectors: the key $\mathbf{k}_t \in \mathbb{R}^d$ and the value $\mathbf{v}_t \in \mathbb{R}^d$, following the nomenclature established in Transformer models. Subsequently, at time $t + \tau$, the value $\mathbf{v}_t$ is accessed using a query vector $\mathbf{q}_{t+\tau} \in \mathbb{R}^d$. This approach mirrors the paradigm of Bidirectional Associative Memories (BAMs) \citep{Kohonen:72,Anderson:72,Nakano:72,Anderson:77}.

The covariance update rule~\citep{Sejnowski:77,Dayan:91} for storing a key-value pair is

\begin{align}
\mathbf{C}_t \ &= \ \mathbf{C}_{t-1} \ + \ \mathbf{v}_{t} \ \mathbf{k}_t^\top \ .
\end{align}

We posit that inputs are first normalized by a layer-norm operation prior to their projection into key and value spaces, ensuring zero mean for these projections. The optimality of the covariance update rule has been established~\citep{Dayan:91} for maximizing the separability of the binary vectors retrieved—a property that corresponds to achieving the greatest possible signal-to-noise ratio. When retrieval is confined to pairwise interactions—resulting in quadratic computational complexity akin to attention mechanisms~\citep{Krotov:16,Krotov:17,Ramsauer:21}—an even greater level of separability may be realized. Notably, the covariance update rule aligns with the formulation of Fast Weight Programmers~\citep{Schmidhuber:92ncfastweights,Schlag:21}, which have since been extended to include a constant decay applied to $\mathbf{C}_{t-1}$ and a fixed learning rate multiplied by 
$\mathbf{v}_{t} \mathbf{k}_t ^\top$~\citep{Ba:16}.
In alignment with this development, we adapt the covariance update mechanism within the LSTM architecture, such that the forget gate performs the role of the decay factor, the input gate serves as the learning rate, and the output gate modulates the retrieved vector.

For this matrix memory, the normalizer state is the weighted sum of key vectors, where each key vector is weighted by the input gate and all future forget gates. Again, the normalizer state keeps record of the strength of the gates. Since the dot product between query and normalizer state can be close to zero, we use the absolute value of this dot product and lower bound it by a threshold (typically 1.0) as done previously~\citep{Sun:23arxiv}. The mLSTM forward pass is:

\begin{align}
\label{eq:mlstm_recurrent_begin}
\mathbf{C}_t \ &= \  f_t \ \mathbf{C}_{t-1} \ + \ 
  i_t \ \mathbf{v}_t \ \mathbf{k}_t^\top & &  &\text{cell state} \\
\mathbf{n}_t \ &= \  f_t \ \mathbf{n}_{t-1} \ + \ 
  i_t \ \mathbf{k}_t & &  &\text{normalizer state} \\
\mathbf{h}_t  \ &= \ \mathbf{o}_t \ \odot \ \tilde{\mathbf{h}}_t
\ , \qquad \qquad 
\tilde{\mathbf{h}}_t \ = \   \mathbf{C}_t \mathbf{q}_t \ / \ 
  \max \left\{ |\mathbf{n}_t^\top \mathbf{q}_t|, 1 \right\} 
   &&&\text{hidden state} \\
\mathbf{q}_t \ &= \ \mathbf{W}_q \ \mathbf{x}_t \ + \ \mathbf{b}_q  & &  &\text{query input} \\
\mathbf{k}_t \ &= \ \frac{1}{\sqrt{d}} \mathbf{W}_k \ \mathbf{x}_t \ + \ \mathbf{b}_k  & &  &\text{key input} \\
\mathbf{v}_t \ &= \ \mathbf{W}_v \ \mathbf{x}_t \ + \ \mathbf{b}_v  & &  &\text{value input} \\
i_t \ &= \ \!\exp \!  \! \left( \tilde{i}_t  \right) \ , 
 \qquad \qquad \ \ \ \ \,
  \tilde{i}_t \ = \ \mathbf{w}^\top_{i} \ \mathbf{x}_t \ + \  b_{i} \ \ \ 
  &&&\text{input gate} \\
f_t \ &= \ \sigma \!  \left(  \tilde{f}_t \right) \ \text{OR} \ \!\exp \!  \! \left(  \tilde{f}_t \right) , \ \ \: \!
\tilde{f}_t \ = \ \mathbf{w}^\top_{f} \ \mathbf{x}_t  \ + \
  b_{f}
  &&& \text{forget gate} \\
\label{eq:mlstm_recurrent_end}
\mathbf{o}_t \ &= \ \sigma \left( \tilde{\mathbf{o}}_t \right) \ , \qquad \qquad \quad \ \ \ \,
  \tilde{\mathbf{o}}_t  \ = \ \mathbf{W}_{\mathbf{o}} \ \mathbf{x}_t \ + \
  \mathbf{b}_{\mathbf{o}}
  &&&\text{output gate} 
\end{align}

The mLSTM architecture allows for several memory cells, similarly to the standard LSTM. Within mLSTM, having multiple heads is functionally the same as employing multiple cells because the design does not incorporate memory mixing. To ensure stable exponential gating behavior in mLSTM, we apply the same stabilization strategies used for sLSTM (refer to Equation~\ref{eq:slstmstabil}). Because memory mixing is absent in mLSTM, its recurrence relations can be adapted to run in parallel. Additional information is provided in Appendix~\ref{sec:appmLSTM}.

\subsection{xLSTM Architecture}
\label{sec:xLSTMarch}

\paragraph{xLSTM Blocks.}
\label{sec:xLSTMblock}
An xLSTM block aims to capture information from the past through a nonlinear transformation into a high-dimensional representation space, facilitating the distinction among various histories or contexts. Distinguishing among different histories is a key condition for accurate prediction of subsequent sequence elements, such as the next token. Our design approach leverages Cover’s Theorem~\citep{Cover:65}, which posits that, through nonlinear mapping, patterns are more likely to become linearly separable in higher-dimensional spaces compared to the original representation. We examine two types of residual block designs: (i) A residual structure featuring up-projection after the nonlinearity (akin to those in Transformers), where historical information is first compressed nonlinearly in the original space before being mapped linearly to a higher-dimensional space, followed by an application of a nonlinear activation, and finally projected linearly back to the original dimension; refer to the left side of Figure~\ref{fig:Backbones} and the third column of Figure~\ref{fig:xlstm_sketch}. 
A comprehensive illustration is provided in Figure~\ref{fig:appsLSTM_detailed} in the appendix. (ii) Alternatively, a residual block with up-projection preceding the nonlinearity (inspired by State Space Models), where the mapping to higher dimensions occurs first,

summarizes previous information in a non-linear manner within a high-dimensional context, followed by a linear transformation projecting back into the input space. For an xLSTM block built with an sLSTM, we utilize the post up-projection component, whereas for one built with an mLSTM, we employ the pre up-projection mechanism, since high-dimensionality increases memory capacity. For a visual explanation, consult the left side of Figure~\ref{fig:Backbones}, the third column in Figure~\ref{fig:xlstm_sketch}, or Figure~\ref{fig:appsLSTM_detailed} located in the appendix.

\paragraph{xLSTM Architecture. }
\label{sec:xLSTMarchitecure}
The xLSTM architecture is created through the residual stacking of modular components~\citep{Srivastava:15,He:16}. Our approach is based on prevalent pre-LayerNorm~\citep{Ba:16b} residual structures, reflecting the architecture favored in state-of-the-art Large Language Models. Refer to the rightmost column in Figure~\ref{fig:xlstm_sketch} for illustration.

\subsection{Memory and Speed Considerations}

In contrast to Transformers, xLSTM architectures maintain linear computational requirements and constant memory complexity relative to sequence length. The compressive nature of xLSTM memory makes it especially favorable for use in industrial contexts and edge device deployments.

Although the memory mechanism in mLSTM operates without additional parameters, it incurs high computational demands via its $d \times d$ matrix memory and corresponding $d \times d$ update operations. Here, increased memory capacity comes at the cost of computational intensity. Despite this, the relevant operations can be executed in parallel on GPUs, resulting in only a slight impact on wall clock performance.

Analogous to approaches such as FlashAttention~\citep{Dao:22,Dao:23} and GLA~\citep{Yang:23arxiv}, mLSTM supports parallelization; in contrast, sLSTM cannot be parallelized owing to its memory mixing involving hidden-to-hidden connections. To address this, we engineered an optimized CUDA implementation with GPU memory management down to the register level, achieving a runtime that is generally under twice as slow as mLSTM.

\section{Related Work}
\label{sec:related}

\paragraph{Linear Attention.} Numerous approaches have been proposed to address the Transformer’s quadratic scaling with respect to context length, aiming to achieve linear complexity in attention mechanisms. Synthesizer generates synthetic attention weights independently of token--token interactions~\citep{Tay:20arxiv}. Linformer employs a low-rank matrix formulation to implement self-attention, enabling linear approximation~\citep{Wang:20arxiv}. Linear Transformer reformulates the attention computation to operate in a linear fashion~\citep{Katharopoulos:20}. Performer utilizes positive orthogonal random features to obtain a linear estimate of the attention softmax~\citep{Choromanski:21}. In Structured Global Convolution (SGConv)~\citep{Li:22} and Hyena Hierarchy~\citep{Poli:23}, attention mechanisms are substituted with efficient long-range convolutional operations.

\paragraph{State Space Models.} State Space Models (SSMs) have garnered significant attention lately, primarily because their computational complexity scales linearly with the context length, and they demonstrate competitive results against Transformer-based architectures. Early developments in this area include the introduction of the Structured State Space sequence model (S4)~\citep{Gu:21}. Subsequent advancements led to the Diagonal State Space (DSS) model~\citep{Gupta:22}, Gated State Space (GSS) models~\citep{Mehta:22}, the S5 model~\citep{Smith:22}, Bidirectional Gated SSM (BiGS)~\citep{Wang:22}, H3 model~\citep{Fu:23}, and Mamba~\citep{Gu:24arxiv}.

\paragraph{Recurrent Neural Networks.} The use of Recurrent Neural Networks (RNNs) has been proposed as an alternative to Transformer and attention mechanisms, primarily due to their efficiency stemming from linear scaling with sequence length. Language modeling experiments utilizing Deep Linear Recurrent Units (LRUs) have yielded encouraging outcomes~\citep{Orvieto:23, De:24arxiv}. Additionally, models like the Hierarchically Gated Linear RNN (HGRN)~\citep{Qin:23} and its successor HGRN2~\citep{Qin:24arxiv} have demonstrated strong performance. Among RNN-based strategies for large language models, RWKV~\citep{Peng:23arxivshort,Peng:24arxivshort} stands out for offering results on par with Transformer architectures.

\paragraph{Gating.}
LSTM's introduction of gating has significantly influenced modern architectures, leading to its re-emergence and novel interpretations in numerous recent models. Variations of the gating concept are incorporated in HGRN~\citep{Qin:23}, HGRN2~\citep{Qin:24arxiv}, Gated Linear Attention (GLA)~\citep{Yang:23arxiv}, Gated State Space (GSS) models~\citep{Mehta:22}, Bidirectional Gated SSM (BiGS)~\citep{Wang:22}, Moving Average Equipped Gated Attention (MEGA)~\citep{Ma:22}, RWKV~\citep{Peng:23arxivshort}, as well as Mamba~\citep{Gu:24arxiv}.

\paragraph{Covariance Update Rule.} 
To bolster storage efficiency, we integrate a matrix memory governed by a covariance update rule within the mLSTM cell. The principle of such update strategies also underpins other techniques, notably Fast Weight Programmers~\citep{Schmidhuber:92ncfastweights,Schlag:21}, RWKV-5 and RWKV-6~\citep{Peng:24arxivshort}, Retention~\citep{Sun:23arxiv}, Linear Transformer~\citep{Katharopoulos:20}, and HGRN2~\citep{Qin:24arxiv}.

\paragraph{Most Related.} The models most similar to xLSTM in principle include Retention~\citep{Sun:23arxiv}, RWKV~\citep{Peng:23arxivshort,Peng:24arxivshort}, and HGRN2~\citep{Qin:24arxiv}. These frameworks employ mechanisms like matrix memory and gating. Nonetheless, unlike the novel sLSTM, these techniques do not incorporate memory mixing. The inclusion of memory mixing addresses challenges related to state tracking, making LSTMs inherently more capable compared to State Space Models (SSMs) and Transformers~\citep{Merrill:24,Deletang:23}. Tasks such as code execution and maintaining coherence of entities across extensive narratives depend on accurate state tracking.

\paragraph{Residually Stacking Architectures.} xLSTM models, in alignment with prevailing trends in large-scale deep learning frameworks, employ residual stacking of architectural modules~\citep{Srivastava:15,He:16}.

This methodology paved the way for advances in deep convolutional networks~\citep{He:16} as well as for the development of Transformers~\citep{Vaswani:17}. Transformers serve as the foundational architecture powering prominent Large Language Models (LLMs), such as GPT-3~\citep{Brown:20short}, ChatGPT~\citep{Schulman:22short}, GPT-4~\citep{Achiam:23gpt4short}, Megatron-LM~\citep{Shoeybi:19}, Gopher~\citep{Rae:21short}, ERNIE 3.0 Titan~\citep{Wang:21arxivshort}, GLaM~\citep{Du:21glamshort}, Chinese M6~\citep{Lin:21}, AlexaTM 20B~\citep{Soltan:22}, OPT~\citep{Zhang:22opt}, Chinchilla~\citep{Hoffmann:22short}, BLOOM~\citep{Scao:22arxivshort}, GLM-130B~\citep{Zeng:22arxivshort}, LaMDA~\citep{Thoppilan:22short}, PaLM~\citep{Chowdhery:22short}, Llama~\citep{Touvron:23arxiv}, Gemini~\citep{GeminiTeam:23arxiv,Reid:24arxiv}.

\section{Experiments}
\label{sec:Exp}

This section presents a comprehensive experimental analysis of xLSTM in comparison to established approaches, with an emphasis on language modeling. Section~\ref{sec:ExpSyntethic} explores xLSTM's performance and properties through a series of synthetic benchmarks. Section~\ref{sec:ExpComparison} details a comparative study of several current language modeling strategies, measuring their validation set perplexity when trained on 15B tokens sourced from SlimPajama~\citep{Soboleva:23}. We additionally execute ablation experiments on xLSTM using this dataset. The scalability of these models is then examined using frameworks similar to those described in \citet{Kaplan:20} and \citet{Brown:20short}.

Section~\ref{sec:ExpLanguage} extends the investigation to a larger-scale language modeling task. Here, xLSTM is evaluated alongside the top-performing models identified in Section~\ref{sec:ExpComparison}, following training on 300B SlimPajama tokens~\citep{Soboleva:23}. This stage of evaluation includes: (1) testing the generalization of these methods to substantially longer input contexts, (2) comparing their validation perplexity and performance on subsequent tasks~\citep{Sutawika:24}, (3) assessing each model on 571 text domains from the PALOMA benchmark~\citep{Magnusson:23arxivshort}, and (4) re-examining scaling characteristics with a dataset twenty times larger than in earlier experiments.

\label{sec:xlstm_blocks_notation}
Throughout all experimental settings, we denote the proportion of mLSTM- to sLSTM-based xLSTM blocks as xLSTM[$a$:$b$], where $a/b$ reflects the specific ratio. As an illustration, xLSTM[7:1] designates a configuration in which seven out of eight blocks are mLSTM-based, with the remaining one being sLSTM-based. In scenarios with a total of 48 blocks—a commonly used size—this convention corresponds to 42 blocks based on mLSTM and 6 based on sLSTM. Additionally, every experiment incorporates up-projection blocks both before and after, applied respectively to mLSTM and sLSTM modules.

\subsection{Synthetic Tasks and Long Range Arena}
\label{sec:ExpSyntethic}
Initially, we evaluate how well the xLSTM's novel exponential gating mechanism with memory mixing operates on formal language tasks~\citep{Deletang:23}. Next, we examine the capabilities of xLSTM's innovative matrix memory by applying it to the Multi-Query Associative Recall task~\citep{Arora:23arxiv}. Lastly, the ability of xLSTM to manage extended sequence processing is measured using benchmarks from the Long Range Arena~\citep{Tay:21}.

\paragraph{Evaluation of xLSTM's Exponential Gating Mechanism with Memory Mixing.} We assess the capabilities of xLSTM's exponential gating combined with memory mixing, a feature designed to address state tracking challenges~\citep{Merrill:24, Merrill:23}. Our evaluation extends the formal language tasks introduced by \citet{Deletang:23}, adapting them to support multi-length training regimes necessary for testing length extrapolation. Comprehensive task descriptions and detailed outcomes are available in Appendix~\ref{sec:appExpSynthetic-formalLang}. In this investigation, xLSTM is benchmarked against several alternative architectures, including Transformers, State Space Models, and Recurrent Neural Networks. Performance is measured in terms of accuracy, specifically on those tokens pertinent to each task, and scaled from 0 (indicative of random guessing) to 1 (representing perfect performance). We compare two-block configurations for multiple models: xLSTM[0:1] (just sLSTM), xLSTM[1:0] (just mLSTM), xLSTM[1:1], Llama, Mamba, RWKV, Retention, Hyena, LSTM, and LSTM integrated into Transformer blocks (LSTM (Block)). Outcomes are presented in Figure~\ref{fig:formal-main}. Architectures such as Transformers and State Space Models, when lacking memory mixing (and therefore state tracking capabilities), are unable to solve tasks requiring regular grammars, such as the parity problem. These findings align with previous research which demonstrates that, without state tracking, Transformers and State Space Models are inherently less expressive than RNNs~\citep{Merrill:24, Merrill:23, Deletang:23}.

\paragraph{Test of xLSTM's Memory Capacities on Associative Recall Tasks.} This study evaluates the new matrix memory mechanism in xLSTM by examining its ability to store and retrieve information in the Multi-Query Associative Recall task~\citep{Arora:23arxiv}. For each input sequence, randomly selected key-value pairs are drawn from an extensive vocabulary, requiring accurate retention for subsequent access. To increase the complexity beyond the original setup, we scale the number of key-value pairs to as many as 256 and extend context lengths up to 2048, thereby providing a more rigorous assessment of model memory capabilities. The experiment compares two-block versions of several architectures: Llama, Mamba, RWKV-5, RWKV-6, xLSTM[1:1], and xLSTM[1:0]. Accuracy in retrieving the key-value pairs serves as the evaluation metric. With memory capacity scaling exponentially in the encoding dimension~\citep{Ramsauer:21}, Transformer models such as Llama are established as the benchmark for this task. Results, as presented in Figure~\ref{fig:mqar-main}, show that xLSTM[1:1] outperforms all non-Transformer models, including for models with fewer parameters. Notably, integrating the sLSTM block does not impair, but actually enhances memory capacity, particularly under the most challenging conditions involving 256 key-value pairs. Further findings are available in Appendix~\ref{sec:appExpSynthetic-mqar}, where extrapolation analyses demonstrate that the augmented memory features of xLSTM permit successful generalization to contexts surpassing the length encountered during training.

\paragraph{Test of xLSTM's Long Context Capabilities on Long Range Arena.} In order to evaluate xLSTM's effectiveness in managing extended sequences and large contextual windows, we benchmark various approaches using the Long Range Arena~\citep{Tay:21}. Across all tasks, xLSTM maintains robust results, indicating that its architecture excels in efficiently tackling diverse challenges associated with long context processing.

For more details, see Appendix~\ref{sec:appExpSynthetic-lra}.

\subsection{Method Comparison and Ablation Study}
\label{sec:ExpComparison}

In order to explore the primary focus of our study—namely, evaluating the performance of our novel LSTM variants at scale within language modeling tasks—we conduct experiments by training xLSTMs, Transformers, State Space Models, among other approaches, utilizing 15B tokens sourced from SlimPajama under identical auto-regressive conditions. The resulting models are evaluated on a common validation set, and for xLSTMs specifically, we carry out ablation experiments to dissect their behavior and contributions.

\paragraph{Assessing xLSTM Against Alternative Approaches.} To facilitate a fair comparison, models were trained on a dataset comprising 15B tokens sourced from SlimPajama~\citep{Soboleva:23}. Their performance was measured using perplexity scores calculated on a separate validation set. The methods evaluated include our proposed xLSTM, alongside GPT-3 (Transformer)~\citep{Brown:20short}, Llama (Transformer)~\citep{Touvron:23arxiv}, H3 (SSM)~\citep{Fu:23}, Mamba (SSM)~\citep{Gu:24arxiv}, RWKV-4 (RNN)~\citep{Peng:23arxivshort}, RWKV-5 (RNN)~\citep{Peng:24arxivshort}, RWKV-6 (RNN)~\citep{Peng:24arxivshort}, GLA (linear Transformer)~\citep{Yang:23arxiv}, HGRN (RNN)~\citep{Qin:23}, HGRN2 (RNN)~\citep{Qin:24arxiv}, RetNet (linear Transformer)~\citep{Sun:23arxiv}, Hyena (linear Transformer)~\citep{Poli:23}, and both xLSTM[1:0] and xLSTM[7:1] variants, as described in Section~\ref{sec:xlstm_blocks_notation}. Training employed mixed-precision computation; however, for RWKV-5, RWKV-6, GLA, and HGRN2, mixed precision was implemented without leveraging PyTorch's automated mixed-precision capabilities (see Appendix Section~\ref{sec:appGenTrainProc} for further information).

The models are divided into three groups: (a) Transformers, (b) State Space Models (SSMs), and (c) Recurrent Neural Networks (RNNs), which are grouped together with linear Transformers. Linear Transformer approaches replace the traditional attention mechanism with linear alternatives. All models were configured to mirror the parameter count of GPT-3 at 350M, utilizing an embedding dimension of 1024 and comprising 24 residual blocks. Among these, only GPT-3 shares its token and output embedding weights, thereby resulting in a slightly reduced parameter count.

According to the results summarized in Table~\ref{tab:spaj15b_model_results}, xLSTM achieves the lowest validation perplexity, outperforming all previously reported methods. Additional information is provided in Appendix~\ref{sec:appExpComparison}. Figure~\ref{fig:temp_sclaw15B} illustrates the scaling dynamics for these experiments, supporting the projection that xLSTM’s advantages persist for larger-scale models.

\paragraph{Ablation Studies.}
As shown in Table~\ref{tab:spaj15b_model_results} and Figure~\ref{fig:temp_sclaw15B}, xLSTM performs strongly on language modeling tasks when trained with 15B tokens from SlimPajama. To systematically investigate the effect of each architectural modification, we incrementally transform a standard LSTM into an xLSTM. The process begins by embedding LSTM layers within pre-LayerNorm residual backbone structures. Subsequently, the architecture is expanded to include a post up-projection block. The final step involves incorporating exponential gating mechanisms and matrix memory.

The results are shown in Table~\ref{tab:ablstudies} (top). 

The results from the ablation studies indicate that the notable performance gains can be credited to both the exponential gating mechanism and the matrix memory. Furthermore, given the pivotal role of gating within RNNs and State Space Models, various gating architectures are examined through ablations. As shown in Table~\ref{tab:ablstudies} (bottom), the findings demonstrate that making each gate both trainable and dependent on the input consistently yields further improvements. Further analyses concerning the specific backbone elements can be found in Appendix~\ref{sec:appAblDetails}.

\subsection{xLSTM as Large Language Model}
\label{sec:ExpLanguage}

To conclude our investigation, we conduct extensive language modeling experiments at scale, evaluating xLSTM's viability as a large language model (LLM). To this end, we expand our training set, utilizing 300B tokens sourced from SlimPajama. This dataset size aligns with that employed by Mamba~\citep{Gu:24arxiv} and Griffin~\citep{De:24arxiv}.

xLSTM is benchmarked against RWKV-4, Llama, and Mamba, which each exemplify a distinct architectural category discussed in Section~\ref{sec:ExpComparison}. RWKV-4 is chosen as the RNN baseline, as only after the commencement of the 300B-token training runs was suitable precision found for RWKV-5, RWKV-6, and HGRN2 (see Appendix~\ref{sec:appGenTrainProc}).
A range of model capacities (125M, 350M, 760M, 1.3B) are trained, all subjected to tests for length extrapolation and validation performance. Their results are further analyzed across downstream task evaluations, language modeling accuracy spanning 571 PALOMA benchmark text domains, and an examination of their scaling law characteristics.

\paragraph{Sequence Length Extrapolation.}
\label{sec:spaj300B_ctx_extrapolation}
To begin, we investigate how well large-scale models with 1.3B parameters—including xLSTM, RWKV-4, Llama, and Mamba—extrapolate to longer sequence lengths. Although these models are initially trained using a context window of 2048 tokens, we assess their performance on sequences up to 16384 tokens. The experimental outcomes are illustrated in Figure~\ref{fig:spaj300B_seq_extrapolation}. Notably, xLSTM consistently exhibits lower perplexity across extended contexts compared to alternative approaches.

\paragraph{Validation Perplexity and Downstream Tasks.}
\label{sec:spaj_downstream_eval}
Next, for every parameter setting, we assess xLSTM, RWKV-4, Llama, and Mamba models using the SlimPajama validation split for next token prediction accuracy, as well as on downstream benchmarks that evaluate abilities in common sense reasoning.

The third column in Table~\ref{tab:lmeval} shows the validation set perplexities achieved by various methods. Among these, xLSTM[1:0] and xLSTM[7:1] consistently yield the lowest validation set perplexity across all model sizes, demonstrating superior performance. The subsequent columns in Table~\ref{tab:lmeval} report results for a range of downstream tasks. Across nearly all tasks and model sizes, xLSTM outperforms other methods; exceptions are seen primarily in the ARC task, where Mamba sometimes achieves the highest scores. Further information is provided in Appendix~\ref{sec:appExpLanguage}.

\paragraph{Performance on PALOMA Language Tasks.} As a third evaluation, we assess the next token prediction accuracy of xLSTM, RWKV-4, Llama, and Mamba across all available model sizes using PALOMA language tasks~\citep{Magnusson:23arxivshort}. The assessment metric is perplexity for next token prediction over a set of 571 text domains, encompassing a wide variety from nytimes.com to r/depression on Reddit. 

Table~\ref{tab:ppleval} presents the token prediction perplexity, separating language modeling tasks (first seven columns) from specialized domain benchmarks (last 5 columns). On these tasks, xLSTM[1:0] outperforms xLSTM[7:1].

In comparison, xLSTM[1:0] achieves a lower perplexity than Mamba in 568 out of 571 domains (99.5\%), outperforms Llama in 486 out of 571 domains (85.1\%), and exceeds RWKV-4 in 570 out of 571 domains (99.8\%). Additional analysis is provided in Appendix~\ref{sec:appExpLanguage}.

\paragraph{Scaling Laws.} In our fourth analysis, we examine the characteristic power-law scaling observed in model performance, which facilitates forecasting outcomes for expanded model sizes~\citep{Kaplan:20,Brown:20short}. As depicted in Figure~\ref{fig:temp_sclaw300B}, the scaling trends are consistent across models, though each exhibits a distinct offset. RWKV-4 demonstrates the lowest performance, while Llama and Mamba outperform it by moderate margins. The xLSTM model surpasses Mamba, showing a comparative improvement similar to the gap between Mamba and Llama. These scaling patterns suggest that as models increase in size, xLSTM will continue to outperform both Transformer and State-Space architectures.

\textbf{Generation Times and Maximal Throughput.} In our final analysis, as shown in Figure~\ref{fig:inference_time_generation}, we evaluate the time required for text generation, while Figure~\ref{fig:inference_time_decoding_throughput} (left) presents the peak throughput achieved by the different xLSTM configurations at the 1.3B parameter scale. We benchmark these results against comparable HuggingFace implementations of Mamba, Llama, and RWKV, where the Llama setup incorporates a static key-value cache. At the moment the experiments were conducted, compilation of the full Transformer Model's cache, as well as applying \texttt{torch.compile} to the Mamba implementation, was not possible.

For our generation experiments, each model is evaluated with a batch size of 1 and a pre-fill setting of 16, the latter parameter chosen specifically to best accommodate Transformer-based models. Figure~\ref{fig:inference_time_generation} illustrates that xLSTM and the other recurrent alternatives—Mamba and RWKV-4—demonstrate linear scaling with respect to sequence length, whereas Llama exhibits quadratic scaling. Throughput during decoding is further analyzed with varying batch sizes and prefill values for the Llama architecture. As depicted in Figure~\ref{fig:inference_time_decoding_throughput} (right), xLSTM supports substantially larger batch sizes in comparison to Llama thanks to its constant memory usage, enabling it to attain superior throughput rates.

\section{Limitations}

(i) Unlike mLSTM, sLSTM's approach to memory mixing prevents parallel processing, resulting in limited parallelization capacity and slower execution. Despite this constraint, we implemented an efficient CUDA kernel for sLSTM, achieving performance within a factor of two compared to our parallel mLSTM kernel. (ii) The present CUDA kernels for mLSTM lack optimization, which leads to runtimes roughly four times longer than FlashAttention or the scan operation in Mamba. Improved CUDA performance could be attained by following optimization strategies similar to FlashAttention. (iii) mLSTM’s matrix memory requires computation over $d \times d$ matrices, increasing computational demands. However, since both memory updating and access are parameter-free and parallelizable via basic matrix algebra, the increased complexity only marginally affects overall execution time. (iv) Careful selection of initialization values for forget gates is essential. (v) As mLSTM’s matrix memory does not depend on sequence length, scaling up sequence length for longer contexts could saturate memory capacity. Nevertheless, this has not yet caused problems for contexts up to 16k tokens, as discussed in Section~\ref{sec:spaj300B_ctx_extrapolation}. (vi) The substantial computational requirements of large-scale language modeling prevented us from thoroughly tuning the architecture or hyperparameters for larger xLSTM variants. We foresee that significant optimization efforts will be necessary for xLSTM models to realize their maximum effectiveness.

\section{Conclusion}
Our investigation has yielded partial insight into the initial inquiry: To what extent can language modeling be advanced by expanding LSTM architectures to billions of parameters? The evidence suggests that such scaling allows performance on par with prevalent frameworks such as Transformers and State Space Models. By augmenting LSTM with exponential gating, memory mixing, and a revised memory mechanism, we propose the xLSTM architecture. Empirical results demonstrate that xLSTM is highly competitive in language modeling tasks, rivaling leading approaches including Transformers and State Space Models. Observations of scaling laws further imply that substantially larger xLSTM configurations may serve as formidable alternatives to established Transformer-based Large Language Models. The xLSTM framework holds significant promise for broader application in areas such as Reinforcement Learning, Time Series Forecasting, and modeling complex physical systems.

\end{document}